\section{A boundary for exact universal padded determinant constructions}

\subsection{O-016: the degree-split exponent is optimal for this exact method}
Fix $b\geq2$, and let $R\geq5$ be the number of variables.
Write
\[
 h=\lfloor b/2\rfloor,\qquad
 D_j(R)=\binom{R+j-1}{j}.
\]
For a Zariski-general homogeneous degree-$b$ form $f$ in
$R$ variables over $\mathbb C$, its exact affine-linear
determinantal complexity satisfies
\[
 \operatorname{dc}(f)\geq
 1+\left\lceil\frac{D_b(R)}{D_h(R)+D_{b-h}(R)}\right\rceil
                         =\Omega_b(R^{\lfloor b/2\rfloor}).
                                                               \tag{O-016.1}
\]
Consequently any construction representing every padded
degree-$b$ form in $r=R+1$ variables by an exact homogeneous
linear determinant needs size
$\Omega_b(r^{\lfloor b/2\rfloor})$. The size
$B(b,r)=O_b(r^{\lfloor b/2\rfloor})$ in O-012.a therefore
has the optimal exponent in $r$ for fixed $b$, within this
exact universal problem.

The known input is Chatterjee--Kumar--Volk,
\emph{Determinants vs. Algebraic Branching Programs}, ITCS~2024,
article~27: Observation~2.6, p.~27:6, and Theorem~3.1,
p.~27:8. They imply that if a homogeneous form has singular
locus of dimension less than $R-4$ and an exact determinant
representation of size $s$, it has a homogeneous algebraic
branching program of width at most $s-1$. Their Definition~2.7,
p.~27:7, counts vertices rather than edges for program size;
Theorem~4.1, p.~27:9, more generally gives size $O(b^5s)$
for all homogeneous forms. The sharper width statement is
the only complexity input needed below. These are known
theorems, not new results of this investigation.

\begin{proof}
There is a nonempty Zariski-open set of degree-$b$ forms
whose projective hypersurface is smooth. To see both openness
and nonemptiness, the incidence defined by all partial
derivatives in the product of the coefficient space with
$\mathbb P^{R-1}$ has closed projection, since the second
factor is projective (EGA~II, Theorem~5.5.3(i), p.~104,
with properness Definition~5.4.1, p.~100), while the Fermat form
$z_1^b+\cdots+z_R^b$ is outside that projection.
For such a form the affine singular locus is exactly the
origin, of dimension zero, which is less than $R-4$.

By the cited regular-representation result, any determinant
representation of such an $f$ of size $s$ gives a homogeneous
branching program of width at most $s-1$. Cut the program
at layer $h$. Its path sum then has the form
\[
                f=\sum_{i=1}^{s-1}P_iQ_i,
       \qquad \deg P_i=h,\quad\deg Q_i=b-h,             \tag{O-016.2}
\]
inserting zero summands if the actual width is smaller.
This is simply the partition of all source-to-sink paths by
their vertex at that layer.

For a fixed integer $s$, the forms in (O-016.2) lie in the
image of the polynomial morphism
\[
 \left(\operatorname{Sym}^h\mathbb C^R
             \times\operatorname{Sym}^{b-h}\mathbb C^R\right)^{s-1}
 \longrightarrow\operatorname{Sym}^b\mathbb C^R,
                   \quad (P_i,Q_i)_i\longmapsto\sum_iP_iQ_i.
\]
The closure of its image has dimension at most
$(s-1)(D_h(R)+D_{b-h}(R))$. If this is strictly smaller
than $D_b(R)$, that image closure is proper. Its union with
the proper closed nonsmooth locus is still proper, since
the ambient affine coefficient space is irreducible.
Outside the union no form has determinant complexity at most
$s$. Taking the largest integer $s$ satisfying the strict
dimension inequality proves (O-016.1).

For fixed $b$, the numerator has degree $b$ in $R$ and
the denominator has degree $\lceil b/2\rceil$, with
positive leading coefficients. This gives the stated
asymptotic exponent. Finally, an exact identity
\[
                    \det M(L,z)=L^{s-b}f(z)
\]
with $M$ homogeneous linear and $f$ independent of $L$
specializes at $L=1$ to an affine-linear determinant for
$f(z)$ of size $s$. A universal padded construction includes
these forms, so it inherits the lower bound. O-012.a has
the matching exponent because
$\lceil(b-1)/2\rceil=\lfloor b/2\rfloor$.
\end{proof}

\paragraph{The base-degree twenty-six threshold.}
Take $r=m^2$ and $b=26$, so $R=m^2-1$. Then
\[
 \frac{D_{26}(R)}{2D_{13}(R)}
   =\frac{13!}{2\cdot26!}\prod_{j=13}^{25}(R+j)
                         =\Theta(m^{26}).              \tag{O-016.3}
\]
For all sufficiently large $m$, this exceeds $m^{25}$.
Thus not every degree-$26$ form in these variables can be
handled by an exact padded determinant of size $m^{25}$.
The uniform second-row constant twenty-five obtained from
an exact universal base-form construction cannot be raised
to twenty-six by merely improving the constant factor in
that universal exact construction. Restricted supports and
permanent-specific families remain available, as already
used in F-015--F-017 and M-011.

\paragraph{Exact complexity versus orbit-closure containment.}
The last conclusion is deliberately about exact determinant
identities. It does not prove that the corresponding padded
linear space is absent from $\Omega_{m^{25}}$, nor that a
multiplicity obstruction exists. A point of an orbit closure
can be a limit of determinants whose specialization at $L=1$
has higher-degree terms that disappear only in the limit.
The homogeneous fixed-degree conversion theorem does not
automatically apply uniformly to those approximating forms.
The source itself discusses the unresolved size-preserving
homogeneous-extraction issue on p.~27:5. No passage from
exact to border determinantal complexity is assumed here.
The generic lower bound is also not an explicit lower bound
for the permanent or for any identified GCT coefficient.

\paragraph{Geometric dictionary.}
The coefficient and product parameter spaces above are affine
schemes, and the dimension argument bounds the dimension of
a scheme-theoretic image after base change to $\mathbb C$.
The smooth-form open set is the complement of the proper
projection of a projective incidence scheme. The padded
matrix identity is an integral morphism, while specialization
at the nonzero section $L=1$ is a fiber operation.
These geometric operations do not identify an image with
its closure or assert lifting through a limiting family.
The distinction is the same coherent-quotient versus explicit
parametrization distinction used throughout the GCT comparison.
